\section{Numerical values behind the figures}
\label{app:data}
\setcounter{table}{0}
\setcounter{figure}{0}

The values plotted in Secs.~\ref{sec:orientation}--\ref{sec:transient} are
collected here so that the figures can be read without the underlying files.
They are printed from the same computation that draws each figure, with the
same model and the same mesh, and they are quoted to two significant digits in
the text. Where a table and a figure disagree in the last digit, the figure is
drawn from the finer grid.

Table~\ref{tab:orientation} is the window-orientation sweep of
Figs.~\ref{fig:orientation} and~\ref{fig:design} at every fifth degree in the
quadrant; the two remaining quadrants follow from the symmetries verified in
Sec.~\ref{sec:orientation}. Table~\ref{tab:thickness} is the thickness sweep of
Fig.~\ref{fig:thickness}, together with the two mesh witnesses that establish
that the reported size effect is a property of the model and not of the
mesh.
Table~\ref{tab:damping} is the damping ratio of the fractional element on the
grid of Fig.~\ref{fig:damping}(b), Table~\ref{tab:convergence} the convergence
study of Fig.~\ref{fig:convergence}, and Table~\ref{tab:ringdown} the decrement
of the traces of Fig.~\ref{fig:ringdown}. Three further tables are collected
here rather than in the body: the closed-form coupling factors of
Sec.~\ref{sec:verification}, the ring-down of the memory study of
Sec.~\ref{sec:memorystudy}, and the calibration protocol of
Sec.~\ref{sec:calib_a_mu}.

% generated by python/make_appendix_tables.py -- do not edit
\begin{table}[!htbp]
\centering
\caption{Window orientation: frequency of the extensional mode and its coupling
factor, for the three aspect ratios of Fig.~\ref{fig:orientation}, on the
\SI{200}{\micro\metre}$\times$\SI{10}{\micro\metre} strip with
$h=\SI{10}{\micro\metre}$, poling $x$, plane stress, $24\times2$ mesh,
$\mu=0$ and $\alpha=1$. Every fifth degree is listed; the figure uses the same
sweep at $5^\circ$ spacing. The remaining orientations follow from
$\theta\to\theta+180^\circ$ and $\theta\to-\theta$, which reproduce the values
below to \num{3e-11} relative (Sec.~\ref{sec:orientation}).}
\label{tab:orientation}
\begin{tabular}{@{}l r r r r r r@{}}
\toprule
& \multicolumn{2}{c}{${\rm AR}=2$} & \multicolumn{2}{c}{${\rm AR}=5$}
& \multicolumn{2}{c}{${\rm AR}=10$} \\
\cmidrule(lr){2-3}\cmidrule(lr){4-5}\cmidrule(lr){6-7}
$\theta$ (deg) & $f_{\rm ext}$ (kHz) & $k^2$ & $f_{\rm ext}$ (kHz) & $k^2$
& $f_{\rm ext}$ (kHz) & $k^2$ \\
\midrule
0 & 5012 & 0.4448 & 4921 & 0.4648 & 4907 & 0.4677 \\
15 & 5031 & 0.4403 & 4947 & 0.4590 & 4934 & 0.4617 \\
30 & 5085 & 0.4281 & 5016 & 0.4433 & 5007 & 0.4455 \\
45 & 5157 & 0.4115 & 5110 & 0.4221 & 5103 & 0.4236 \\
60 & 5227 & 0.3950 & 5200 & 0.4009 & 5196 & 0.4018 \\
75 & 5278 & 0.3830 & 5265 & 0.3856 & 5263 & 0.3859 \\
90 & 5296 & 0.3786 & 5289 & 0.3799 & 5288 & 0.3801 \\
\bottomrule
\end{tabular}
\end{table}

% generated by python/make_appendix_tables.py -- do not edit
\begin{table}[!htbp]
\centering
\caption{Thickness sweep behind Fig.~\ref{fig:thickness}:
$L=\SI{200}{\micro\metre}$, window half-widths
$3\times1$\,\si{\micro\metre}, $\mu=10^{-6}$\,\si{\coulomb\per\metre} for the
bending columns and $\mu=0$ for the extensional ones, $\alpha=1$. The
$24\times2$ and $96\times4$ columns are the two mesh witnesses; the
$96\times4$ values were computed at the two ends of the range only. The
fundamental frequency is that of the same $48\times4$ model.}
\label{tab:thickness}
\begin{tabular}{@{}r r r r r r r@{}}
\toprule
& & \multicolumn{3}{c}{bending, $k^2_{\rm flexo}$}
& \multicolumn{2}{c}{extensional, $k^2_{\rm piezo}$} \\
\cmidrule(lr){3-5}\cmidrule(lr){6-7}
$h$ (\si{\micro\metre}) & $f_1$ (kHz) & $48\times4$ & $24\times2$ & $96\times4$
& $48\times4$ & $24\times2$ \\
\midrule
10 & 210.4 & 1.744e-04 & 1.25e-04 & 1.7444e-04 & 0.55814 & 0.55797 \\
7 & 148.7 & 3.510e-04 & 2.58e-04 & -- & 0.55831 & 0.55811 \\
5 & 108.0 & 6.587e-04 & 5.08e-04 & -- & 0.55843 & 0.55819 \\
3.5 & 78.2 & 1.218e-03 & 9.95e-04 & -- & 0.55851 & 0.55824 \\
2.7 & 62.9 & 1.808e-03 & 1.53e-03 & -- & 0.55855 & 0.55826 \\
2 & 50.2 & 2.642e-03 & 2.32e-03 & -- & 0.55858 & 0.55828 \\
1.4 & 40.3 & 3.608e-03 & 3.26e-03 & -- & 0.55860 & 0.55829 \\
1 & 34.7 & 4.174e-03 & 3.85e-03 & 4.1589e-03 & 0.55861 & 0.55829 \\
\bottomrule
\end{tabular}
\end{table}

% generated by python/make_appendix_tables.py -- do not edit
\begin{table}[!htbp]
\centering
\caption{Damping ratio $\zeta=-\mathrm{Re}\,s/(\omega_0|s|)$ of the fractional
Kelvin--Voigt element of \eqref{eq:kvfreq}, on the grid behind
Fig.~\ref{fig:damping}(b), from the complex pole with branch continuation. The
ratio rises with both $\alpha$ and $\omega\tau$ and reaches unity at
$\alpha=1$, $\omega\tau=10$, where the integer-order element is critically
damped. The band used for the strip mode of this paper is
$\omega\tau\le1$.}
\label{tab:damping}
\begin{tabular}{@{}l r r r r r@{}}
\toprule
& \multicolumn{5}{c}{$\omega\tau$} \\
\cmidrule(lr){2-6}
$\alpha$ & 0.1 & 1 & 3 & 10 & 30 \\
\midrule
0.4 & 0.097 & 0.188 & 0.237 & 0.286 & 0.320 \\
0.5 & 0.100 & 0.245 & 0.331 & 0.408 & 0.452 \\
0.6 & 0.095 & 0.307 & 0.446 & 0.551 & 0.594 \\
0.7 & 0.086 & 0.371 & 0.585 & 0.704 & 0.735 \\
0.8 & 0.075 & 0.432 & 0.749 & 0.848 & 0.862 \\
0.9 & 0.062 & 0.480 & 0.872 & 0.939 & 0.941 \\
1.0 & 0.050 & 0.500 & 1.000 & 1.000 & 1.000 \\
\bottomrule
\end{tabular}
\end{table}

% generated by python/make_appendix_tables.py -- do not edit
\begin{table}[t]
\centering
\caption{Convergence of the extensional mode behind
Fig.~\ref{fig:convergence}: plate \SI{200}{\micro\metre}$\times$\SI{10}{\micro\metre},
poling $x$, plane stress, $\mu=0$, $a_1\times a_2=3\times1$\,\si{\micro\metre}.
The coupled frequency and the mode-based coupling are those of the
eigenvalue problem; the errors are referred to the $48\times4$ mesh.}
\label{tab:convergence}
\begin{tabular}{@{}l r r r r r@{}}
\toprule
mesh & dof & $f$ (kHz) & $k^2$ & err $f$ & err $k^2$ \\
\midrule
$12\times2$ & 312 & 0.87 & 1.0000 & 1.60e-03 & 8.3e-11 \\
$24\times2$ & 600 & 0.87 & 1.0000 & 1.86e-04 & 1.2e-11 \\
$36\times2$ & 888 & 0.87 & 1.0000 & 1.83e-05 & 1.1e-12 \\
$48\times4$ & 1960 & 0.87 & 1.0000 & 0.00e+00 & 0.0e+00 \\
\bottomrule
\end{tabular}
\end{table}

% generated by python/make_appendix_tables.py -- do not edit
\begin{table}[!htbp]
\centering
\caption{Ring-down behind Fig.~\ref{fig:ringdown}: the decrement $\delta$ of
the tip displacement over the first two peaks, converted with the exact
relation $\zeta=\delta/\sqrt{4\pi^{2}+\delta^{2}}$ of \eqref{eq:zeta}, for
each Caputo order. The entry is a damping ratio only at $\alpha=1$; for
$\alpha<1$ it is the first-period decrement, because the free decay of a
fractional element is not exponential. The mode is the extensional one at \SI{6727.3}{kHz}
($k^2=0.4583$), with $\tau=\SI{3}{\nano\second}$ and
$\omega\tau=0.13$; the open-circuit voltage at the initial instant is
\SI{139.07}{\volt} for a \SI{10}{\nano\metre} tip displacement. The decrement is
not monotonic in $\alpha$: it peaks near $\alpha=0.6$ and falls again as the
response becomes slower and more heavily damped.}
\label{tab:ringdown}
\begin{tabular}{@{}l r@{}}
\toprule
$\alpha$ & $\zeta$ of the first period \\
\midrule
1.0 & 0.06343 \\
0.95 & 0.07011 \\
0.9 & 0.07688 \\
0.8 & 0.08988 \\
0.7 & 0.10008 \\
0.6 & 0.10408 \\
0.5 & 0.09890 \\
0.4 & 0.08461 \\
\bottomrule
\end{tabular}
\end{table}


\FloatBarrier

\begin{table}[!htbp]
\centering
\small
\caption{Normalized modulus magnitude $|K^{*}|/K$ at matched $\tau$ for the
integer-order Kelvin--Voigt element and for the fractional Zener element whose
relaxation-strength ratio, $r=1.1$, is fitted to the spring-pot at
$\omega\tau=0.3$ and then held fixed, at two Caputo orders.}
\label{tab:modulus}
\begin{tabular}{@{}lcccc@{}}
\toprule
$\omega\tau$ & \multicolumn{2}{c}{$\alpha=1$} & \multicolumn{2}{c}{$\alpha=0.7$} \\
\cmidrule(lr){2-3}\cmidrule(lr){4-5}
 & KV & Zener ($r=1.1$) & KV & Zener ($r=1.1$) \\
\midrule
0.3  & $1.0440$ & $1.0086$ & $1.2555$ & $1.0244$ \\
1    & $1.4142$ & $1.0512$ & $1.7053$ & $1.0504$ \\
10   & $10.050$ & $1.0991$ & $5.5380$ & $1.0894$ \\
100  & $100.005$ & $1.1000$ & $25.588$ & $1.0981$ \\
\bottomrule
\end{tabular}
\end{table}

\FloatBarrier

%% ---- tables moved out of the text in the final revision ----

\begin{table}[!htbp]
\centering
\small
\caption{The coupling factor obtained from the constitutive tensors alone, for
the four kinematics, together with the coupling strength
$e^{\mathsf T}C^{-1}e/\kappa$ from which it follows. The record is that of
Table~\ref{tab:params}, in plane stress and in plane strain. The three finite
element rows below the rule are all mode-free values of \eqref{eq:kstatic},
plane stress, poling $x$, on the meshes named, with the window switched off;
Table~\ref{tab:coupling} sets them beside the same measure with the window on.}
\label{tab:closedform}
% table generated from the archived results
\begin{tabular}{@{}lcc@{}}
\toprule
& $x$ & $k^2$ \\
\midrule
3-D free bar & 1.3647 & 0.5771 \\
3-D uniaxial ($k_{33}^2$) & 0.5996 & 0.3748 \\
2-D plane stress, from the record & 1.2768 & \textbf{0.5608} \\
\midrule
FEM, plane stress, free, $48\times4$ & 1.2768 & \textbf{0.56079} \\
FEM, plane stress, clamped, $24\times2$ & 1.2668 & 0.55884 \\
FEM, plane stress, clamped, $48\times4$ & 1.2705 & 0.55957 \\
FEM, plane stress, clamped, $96\times4$ & 1.2723 & 0.55992 \\
\midrule
2-D plane strain, from the record & 1.0541 & 0.5132 \\
FEM, plane strain, clamped, $48 \times 4$ & 1.0471 & 0.5115 \\
\bottomrule
\end{tabular}

\end{table}

\begin{table}[!htbp]
\centering
\small
\caption{Ring-down of the extensional mode (mode 4, \SI{6.727}{MHz},
$k^2_{4}(K_{uu})=0.4650$, $\omega\tau=0.3$,
$\tau=\SI{7.0974}{\nano\second}$), for memory carried on the condensed
stiffness $K_{\rm eff}$ and on the short-circuit stiffness $K_{uu}$. At
$\alpha=1$ the entry is the damping ratio \eqref{eq:zeta} of the exponential
decay. At $\alpha=0.7$ the free decay of a fractional element is not
exponential, so no single damping ratio exists: the entry there is the
first-period local ratio of the first two peaks, and the successive decrements
are listed in the text. The last row is the ratio of the two entries.}
\label{tab:memory}
\begin{tabular}{@{}lcc@{}}
\toprule
memory matrix & $\alpha=1$: $\zeta_{\rm eff}$ &
  $\alpha=0.7$: first-period $\zeta$ \\
\midrule
$K_{\rm eff}$ & $0.15003$ & $0.1942$ \\
$K_{uu}$      & $0.06746$ & $0.0831$ \\
\midrule
ratio $K_{uu}/K_{\rm eff}$ & $0.450$ & $0.428$ \\
\bottomrule
\end{tabular}
\end{table}

\begin{table}[!htbp]
\centering
\small
\caption{Calibration protocol. Each step fixes one quantity, and the steps are
ordered so that the later ones check the earlier ones. The last row is the test
of the second-moment description itself, which this paper does not attempt.}
\label{tab:protocol}
\begin{tabular}{@{}p{0.03\textwidth}p{0.30\textwidth}p{0.28\textwidth}p{0.26\textwidth}@{}}
\toprule
\# & measurement & observable & quantity fixed \\
\midrule
1 & bending of strips cut along $x$ and $y$, three or more thicknesses in
$h/\sqrt{M_{22}}\in[2,20]$ & $D_{\rm app}/D_{\rm cl}$ against $h$ &
$M_{22}$ through the thickness \\
2 & the same strips cut at $0^\circ,15^\circ,\dots,90^\circ$ &
$D_{\rm app}(\theta)$ and $f_{1}(\theta)$ & $M_{11}/M_{22}$ and $\theta$ \\
3 & ultrasonic dispersion & $c(k)$ over the band where the gradient term is
visible & $(1/4)M_{mn}$, without a boundary condition \\
4 & boundary-layer profile by digital image correlation & decay length of the
strain at a clamped edge & $M_{22}$ from one specimen \\
5 & image of the microstructure & $M_{mn}=\langle s_{m}s_{n}\rangle_{\mathcal
W}$ by numerical integration & the anisotropy of the window \\
6 & specimens spanning the aspect ratio and the window angle & $k^{2}$ and
$f_{1}$ against ${\rm AR}$ and $\theta$ & a prediction, not a fit \\
7 & reconstructed microstructure & decay lengths of the two-point correlation
& adequacy of the second-moment description \\
\bottomrule
\end{tabular}
\end{table}

\FloatBarrier
